\section{The Third Week of Advent: Rejoice in the Lord}\label{sec:week3}

At the midpoint of Advent, the violet lightens to rose. From the first Latin word of its ancient Entrance Antiphon---taken from Philippians~4:4--5, \textit{Gaudete in Domino semper} (``Rejoice in the Lord always'')---the Third Sunday of Advent has long been called \textit{Gaudete} Sunday. The Church lights the rose candle on the wreath and permits rose vestments at the altar (GIRM, no.~346), not because watchfulness is over, but because the One for whom we watch is already standing at the gate.

\subsection{The Third Sunday of Advent (\textit{Gaudete}) Across Cycles A, B, and C}

\begin{threestable}{Cycle}{Appointed Readings}{Thematic Arc}
Year A & Isaiah~35:1--6a, 10; Psalm~146 (145); James~5:7--10; Matthew~11:2--11 & The desert blooming and the blind seeing; the farmer waiting patiently for rain; John's question from prison answered by Christ's saving works. \\
\addlinespace
Year B & Isaiah~61:1--2a, 10--11; Canticle Luke~1:46--48, 49--50, 53--54 (\textit{Magnificat}); 1~Thessalonians~5:16--24; John~1:6--8, 19--28 & Anointed to bring glad tidings to the poor; Mary's \textit{Magnificat} sung as the Responsorial Canticle; ``Always rejoice. Pray without ceasing''; John confessing ``I am not the Christ.'' \\
\addlinespace
Year C & Zephaniah~3:14--18a; Canticle Isaiah~12:2--3, 4bcd, 5--6; Philippians~4:4--7; Luke~3:10--18 & The Lord dancing in joy over Sion; drawing waters with joy from the Savior's fountains; ``Rejoice in the Lord always \dots\ The Lord is nigh''; concrete charity in answer to ``What then shall we do?'' \\
\end{threestable}

\subsubsection{Year A: Courage for the Fainthearted and Patience in Waiting}

Isaiah~35:1--6a, 10 opens Year~A's \textit{Gaudete} Sunday with a landscape transformed by God's nearness: ``The land that was desolate and impassable shall be glad, and the wilderness shall rejoice, and shall flourish like the lily. \dots\ Say to the fainthearted: Take courage, and fear not: behold your God will bring the revenge of recompense: God himself will come and will save you.'' Psalm~146 (Vulgate~145) names the exact works of that coming God---giving food to the hungry, freeing the fettered, enlightening the blind, supporting the fatherless and the widow---and St.~James offers the homely patience of agriculture (James~5:7--8): ``Be patient therefore, brethren, until the coming of the Lord. Behold, the husbandman waiteth for the precious fruit of the earth: patiently bearing till he receive the early and latter rain. Be you therefore also patient and strengthen your hearts: for the coming of the Lord is at hand.''

In Matthew~11:2--11, John the Baptist sits in Herod's prison and sends two disciples to ask, ``Art thou he that art to come, or look we for another?'' Jesus does not answer with an abstract claim; He points to Isaiah~35 happening before their eyes---the blind see, the lame walk, lepers are cleansed, the deaf hear, the dead rise, and the poor have the gospel preached to them---and adds a beatitude that searches every Advent heart: ``And blessed is he that shall not be scandalized in me.'' Preaching on this passage (\textit{Homily~6 on the Gospels}, read at Advent Mattins), Pope St.~Gregory the Great asks how Christ's miracles could ever be a stumbling-block, and points straight to Calvary: men are tempted to despise the humble suffering of the Author of life, yet ``the more humbling God hath undergone for man's sake, the more worthy is He that man should worship Him.''

\subsubsection{Year B: The Voice Who Points Away From Himself}

Year~B pairs Isaiah~61:1--2a, 10--11---``The spirit of the Lord is upon me, because the Lord hath anointed me: he hath sent me to preach to the meek, to heal the contrite of heart''---with the rare liturgical gift of singing Mary's \textit{Magnificat} (Luke~1:46--54) directly as the Responsorial Canticle of the Mass. St.~Paul's brief, luminous exhortation to the Thessalonians follows (1~Thessalonians~5:16--18, 23--24): ``Always rejoice. Pray without ceasing. In all things give thanks for this is the will of God in Christ Jesus concerning you all. \dots\ He is faithful who hath called you, who also will do it.''

John~1:6--8, 19--28 presents the Baptist under interrogation by priests and Levites from Jerusalem. Asked ``Who art thou?'' he strips himself of every borrowed glory: ``And he confessed and did not deny: and he confessed: I am not the Christ.'' He claims only to be ``the voice of one crying in the wilderness,'' pointing to One already standing unrecognized in their midst. Both St.~Gregory the Great (\textit{Homily~7 on the Gospels}) and St.~Augustine (\textit{Tractates on John} IV.1, IV.7) marvel at John's transparent humility. Gregory writes that ``in that he arrogated not to himself the name of Christ, he became a member of Christ.'' Augustine contrasts John with every earthly court crier:

\begin{quote}
Would it not have seemed to you that a herald would have cried, ``Go away, make room.'' Instead of the herald's cry ``Go away,'' John says ``Come.'' The herald makes men stand back from the judge; to the Judge John calls. Yes, indeed, John calls men to the lowly One, that they may not experience what He will be as the exalted Judge.
\end{quote}

\subsubsection{Year C: The God Who Rejoices Over His People}

Year~C lifts \textit{Gaudete} joy to astonishing tenderness. Zephaniah~3:14--18a tells the daughter of Sion not merely to rejoice in God, but to realize that God Himself rejoices over her: ``The Lord thy God in the midst of thee is mighty, he will save: he will rejoice over thee with gladness, he will be silent in his love, he will be joyful over thee in praise.'' The Responsorial Canticle from Isaiah~12:2--6 answers---``Thou shalt draw waters with joy out of the saviour's fountains''---and Philippians~4:4--7 gives \textit{Gaudete} Sunday its signature charter:

\begin{quote}
Rejoice in the Lord always: again, I say, rejoice. Let your modesty be known to all men. The Lord is nigh. Be nothing solicitous: but in every thing, by prayer and supplication, with thanksgiving, let your petitions be made known to God. And the peace of God, which surpasseth all understanding, keep your hearts and minds in Christ Jesus.
\end{quote}

Luke~3:10--18 immediately tests whether that joy is real by putting practical questions in the mouths of ordinary people---crowds, tax collectors, and soldiers---asking John, ``What then shall we do?'' John does not tell them to abandon their stations in life; he tells the one with two coats to share with the one who has none, the tax collector to take nothing beyond what is appointed, and the soldier to practice neither extortion nor false accusation. Christian joy is never severed from justice and mercy.

\subsection{St.~Leo's December Fast and the Postconciliar Fate of the Ember Days}

In the 1962 Roman Rite, the Wednesday, Friday, and Saturday after the Third Sunday of Advent are the \textbf{Winter Ember Days} (\textit{Quatuor Tempora Adventus}), anchored by ancient stational Masses at St.~Mary Major, the Twelve Holy Apostles, and St.~Peter's, and accompanied by a rich treasury of fifth-century sermons by Pope St.~Leo the Great.

In the postconciliar reform of 1969, the four universal Ember weeks were removed from the General Roman Calendar. Under \textit{Universal Norms on the Liturgical Year and the Calendar} (nos.~45--47), the Holy See entrusted seasonal days of prayer for human needs, the fruits of the earth, and human labor to the pastoral decision of each episcopal conference, to be adapted to local climate and circumstance. In the United States, the bishops have designated Thanksgiving Day as a national day of thanksgiving and allow local dioceses to keep traditional or seasonal days of prayer without imposing universal Advent Ember fasts.

Yet the spiritual marrow of St.~Leo's December preaching remains vital for every postconciliar Christian during the Third Week. In his sermons on the tenth-month fast---still prized in the Church's patristic treasury---Leo joined prayer, self-restraint, and almsgiving into a single living cord (\textit{Sermon~12.4}): ``For by prayer we seek to propitiate God, by fasting we extinguish the lusts of the flesh, by alms we redeem our sins \dots\ Because in prayer faith remains stedfast, in fastings life remains innocent, in almsgiving the mind remains kind.'' In his Second Sermon on the December Fast (read in the traditional Advent Office; \textit{Sermon~13}), Leo coined a line no Christian should ever forget: ``Let us spend in good deeds what we take from indulgence. Let our fast become the banquet of the poor.'' And in \textit{Sermon~19.2--3}, he warned that bodily abstinence means little unless ``the mind fast from vices,'' urging believers to rejoice ``that by you the Lord feeds and clothes His poor.'' Even without a canonical Ember precept, setting aside Wednesday and Friday of the Third Week for simpler meals and deliberate charity keeps Rome's oldest December wisdom alive inside the postconciliar year.

\subsection{The Variable Weekdays of the Third Week}

Here appears one of the most distinctive calendar mechanics of the postconciliar Lectionary (\textit{Ordo Lectionum Missae}, nos.~187--191, followed by the Advent weekday Gospel Acclamation verses at no.~192). Because the second phase of Advent always begins on \textbf{December~17}, the weekday course of Week~III prints only \textbf{five} positions---Monday through Friday---and \textbf{no Saturday of Week~III}, because Saturday of Week~III can never fall earlier than December~17 (when First Sunday is November~27, Third Sunday is December~11, Week~III Monday--Friday occupy December~12--16, and Saturday is December~17).

Depending on the calendar date of the First Sunday of Advent, between \textbf{zero and five} of these Week~III weekdays actually occur in any given year:

\begin{itemize}
  \item When First Sunday falls on \textbf{November~27} (e.g., 2033), all five Week~III weekdays occur (December~12--16, with December~12 taken by Our Lady of Guadalupe in the United States).
  \item When First Sunday falls on \textbf{November~29} (e.g., 2026), three Week~III weekdays occur: Monday~December~14, Tuesday~December~15, and Wednesday~December~16, before December~17 begins on Thursday.
  \item When First Sunday falls on \textbf{December~2 or December~3} (e.g., 2028, 2029, 2034, 2035), \textbf{zero} Week~III weekdays occur at Mass, because Monday after the Third Sunday is already December~17 or December~18.
\end{itemize}

Whenever they occur, the five Week~III weekdays move from Balaam's star to the Baptist's lamp:

\begin{itemize}
  \item \textbf{Monday:} Numbers~24:2--7, 15--17a (``A star shall rise out of Jacob and a sceptre shall spring up from Israel'') and Matthew~21:23--27, the chief priests questioning Jesus' authority and refusing to answer where John's baptism came from.
  \item \textbf{Tuesday:} Zephaniah~3:1--2, 9--13 (God leaving in Israel's midst ``a poor and needy people'' who trust in the name of the Lord) and Matthew~21:28--32, the two sons sent into the vineyard, showing that repentant sinners enter the kingdom ahead of self-satisfied talkers.
  \item \textbf{Wednesday:} Isaiah~45:6b--8, 18, 21b--25 (``Drop down dew, ye heavens, from above, and let the clouds rain the just: let the earth be opened, and bud forth a saviour'') and Luke~7:18b--23, Luke's parallel of John's inquiry from prison.
  \item \textbf{Thursday:} Isaiah~54:1--10 (the Lord's everlasting covenant of peace to the barren, forsaken city: ``the mountains shall be moved \dots\ but my mercy shall not depart from thee'') and Luke~7:24--30, Jesus praising John the Baptist as more than a prophet.
  \item \textbf{Friday:} Isaiah~56:1--3a, 6--8 (``My house shall be called the house of prayer, for all nations'') and John~5:33--36, John remembered as ``a burning and a shining light,'' yet Christ's own works bearing a still greater testimony.
\end{itemize}

In years when December~17 arrives early and suppresses some or all of these Week~III weekdays at Mass, a reader keeping the daily rule can read Isaiah~45:8 (\textit{Rorate caeli}) and the skipped Week~III readings quietly on the evening of December~16, letting their longing flow directly into the Great O~Antiphons.
